\chapwithtoc{Introduction}

% TODO hodne vypichnout motivaci

% Code optimization has been the bane of existence of many programmers\dots And hardware makes it even more important\dots

% There are many aspects to code optimization\dots (algorithmic, computational, memory, parallelism, etc.) and many levels at which it can be performed\dots (source code, intermediate representation, machine code, hardware, etc.)

% There are many formal approaches to code optimization\dots (profile-guided optimization, autotuning, polyhedral model, etc.)

% However, these approaches offer limited help with optimizing data layouts\dots (proving semantic equivalence is hard, search space is huge, etc.); so a human programmer is often needed\dots

% In this thesis, we explore two approaches to help with data layout optimization: layout-agnostic abstractions and LLM-directed code optimization that can complement each other\dots

% \begin{XXX}
% We should at least imply some research questions here\dots
% \end{XXX}

Code optimization is a \xxx{bane} of existence of many software developers, especially those working on high-performance computing applications. Our particular focus lies within the domain of scientific computing \xxx{and general numerical computing}, where performance optimizations can bring orders of magnitude improvements in execution time, runtime resources utilization, and energy efficiency.
\xxx{With the advent of Large Language Models (LLMs) and their increasing use, code optimization of numerical computing applications is becoming even more important for their efficient execution.}

\xxx{These} computations run on a wide variety of hardware architectures, including clusters of multi-core CPUs and GPUs. Such systems are complex in terms of memory hierarchies, communication bandwidths and latencies, parallelism capabilities, and many other factors that can affect the performance of applications. Finding the right optimizations for a given application and hardware architecture is a challenging task; our work is particularly focused on optimizing data layouts and traversal orders, which alone can have a significant impact on performance. Consider a simple example of a matrix multiplication operation in \cref{lst:matrix-multiplication}, which is a fundamental building block in many scientific computing applications.
When running this code on a modern hardware, the performance of the algorithm is \xxx{degraded} due to iterating over whole rows and columns of the input matrices in the innermost loop, which results in inefficient use of the memory hierarchy. The rows and columns of the input matrices do not fit within the cache of the running hardware, so they have to be loaded repeatedly from the main memory of the machine.

\begin{lstlisting}[language=C++, caption={Matrix multiplication}, label={lst:matrix-multiplication}]
float C[M][N], A[M][K], B[K][N];
for (int i = 0; i < M; ++i) {
    for (int j = 0; j < N; ++j) {
        C[i][j] = 0;
        for (int k = 0; k < K; ++k) {
            C[i][j] += A[i][k] * B[k][j];
        }
    }
}
\end{lstlisting}

A simple optimization to this code is to change split each of the three loops into two nested loops, and reorder them so that the innermost loop iterates over small blocks of the input matrices that fit within the cache of the running hardware. This optimization is known as tiling or blocking, and by itself can bring significant performance improvements. A different optimization technique is parallelization, which can be achieved by recognizing that the computation of each \lstinline{C[i][j]} is independent of the others --- on a system equipped with Graphical Processing Units (GPUs) with thousands of cores, we can compute each \lstinline{C[i][j]} in a separate thread, which can lead to orders of magnitude improvements in performance.

Ideally, we would like to be able to experiment with many different optimizations combining various data layout and traversal order transformations and parallelization strategies, and find the best combination for a given application and hardware architecture. We \xxx{have achieved greatness} by developing flexible abstractions for data layout and traversals that allow developers to experiment with the aforementioned optimizations without needing to change the core algorithmic logic, and by exploring the potential of LLMs to assist with code optimization \xxx{some continuation here\dots}.

% TODO: odkryt vraha uz tady
